% F082｜06 总结
\clearpage
\begin{summarybox}
\textbf{\textcolor{CSummary}{一句话理解}}

\textbf{两个根号是到同一个动点的两段距离；相减看相差多少。差的绝对值不超过两个固定点的距离，但能不能等于，还要看延长线是否真的碰到 $x$ 轴。}

\medskip
\textbf{\textcolor{CSummary}{一眼识别}}
\[
 f(x)=\sqrt{(x-m)^2+a^2}-\sqrt{(x-n)^2+b^2},
 \quad a,b>0,\ x\in\mathbb R.
\]
两个固定点取 $A(m,a)$、$B(n,b)$，动点取 $P(x,0)$，则 $f(x)=PA-PB$。

\textbf{\textcolor{CSummary}{必须记清的两个数}}
\[
 \boxed{d=|m-n|},\qquad
 \boxed{D=\sqrt{(m-n)^2+(a-b)^2}}.
\]

\textbf{\textcolor{CSummary}{完整值域速查}}
\[
 \boxed{\begin{array}{c|c}
 a>b&(-d,D]\\
 a<b&[-D,d)\\
 a=b,\ m\ne n&(-d,d)\\
 a=b,\ m=n&\{0\}
 \end{array}}
\]

\textbf{\textcolor{CSummary}{极值发生在哪里？}}

当 $a\ne b$ 时：
\[
 \boxed{x_0=\frac{bm-an}{b-a}}.
\]
$a>b$ 在这里\textbf{取最大值 $D$}；$a<b$ 在这里\textbf{取最小值 $-D$}。$a=b$ 不能套此式。

\end{summarybox}
\clearpage
\begin{summarybox}
\textbf{\textcolor{CSummary}{三种证明怎么选？}}
\begin{description}[style=nextline,leftmargin=2.4em,labelsep=.5em,itemsep=.25em]
\item[\textbf{最易懂：}] 初中几何法：先用两点间线段最短限制距离差，再用延长线找等号。
\item[\textbf{最快算：}] 有理化与平方：用 $(bt-as)^2\geq0$ 求 $D$；把高度调成一样，判断另一边界一定取不到。
\item[\textbf{最常规：}] 高中导数法：确定先增后减还是先减后增，特别适合有 $x$ 范围限制的题。
\end{description}

\textbf{\textcolor{CSummary}{一条高考速答技巧}}

根号一减一，有理化后\textbf{分母恒正、分子是一次式}：
\[
 f(x)=\frac{(n-m)(2x-m-n)+a^2-b^2}{
 \sqrt{(x-m)^2+a^2}+\sqrt{(x-n)^2+b^2}}.
\]
要判断函数正负、求零点，可以\textbf{只看分子}，不必求导。

\textbf{最后一眼检查：}题目中的 $x$ 是所有实数吗？$a,b$ 哪个大？边界到底\textbf{能取到}还是\textbf{只能靠近}？三个问题都答完，括号才不会写错。
\end{summarybox}
